\section{Comandos básicos do Terminal}

\begin{frame}{Por que aprender comandos do terminal?}
  \begin{itemize}
    \item O Git é otimizado para uso via linha de comando; entender o terminal torna operações mais rápidas e repetíveis.
    \item Automatização de tarefas e scripts dependem do terminal.
  \end{itemize}
\end{frame}

% Ensinar comando diff e patch, hash (como o Git verifica se corrompeu)

\begin{frame}[fragile]{Comandos essenciais}
  \begin{itemize}
    \item \texttt{pwd} — mostra o diretório atual
    \item \texttt{ls} — lista arquivos (\texttt{ls -la} mostra detalhes)
    \item \texttt{cd} — muda de diretório
    \item \texttt{mkdir}, \texttt{rmdir} — criar/ remover diretórios
    \item \texttt{rm}, \texttt{cp}, \texttt{mv} — manipular arquivos
    \item \texttt{cat}, \texttt{less}, \texttt{head}, \texttt{tail} — visualizar conteúdo
    \item \texttt{touch} — criar arquivo vazio ou atualizar timestamp
    \item \texttt{grep} — buscar texto em arquivos
    \item \texttt{history} e \texttt{!! / !n} — repetir comandos
  \end{itemize}
\end{frame}
